\documentclass[aspectratio=169,11pt]{beamer}

\usetheme{Madrid}
\usecolortheme{default}
\usefonttheme{professionalfonts}
\setbeamertemplate{navigation symbols}{}
\setbeamertemplate{footline}[frame number]

\usepackage{amsmath, amssymb, booktabs, array}

\title{Violence, Safety, Mobility, And Labor-Market Access}
\subtitle{Gender Study --- Week 7}
\author{Teaching deck draft}
\date{}

\begin{document}

\begin{frame}
  \titlepage
\end{frame}

\begin{frame}{Central question}
\begin{itemize}
    \item How do safety, harassment, violence, mobility constraints, and reproductive autonomy shape labor-market access?
    \item Which margins move: participation, search, acceptance, occupations, quits, retention, wages, or welfare?
    \item What data make hidden harms visible?
    \item When are wages and employment incomplete welfare measures?
\end{itemize}
\end{frame}

\begin{frame}{Why Week 7 is distinct from Week 6}
\begin{tabular}{p{0.35\linewidth} p{0.50\linewidth}}
\toprule
Week 6 & Week 7 \\
\midrule
Policy incidence & Constraints and hidden harms \\
Which policy margin moves? & Which jobs are actually feasible? \\
Who bears incidence? & Who faces risk, control, or restricted mobility? \\
Firm and household adjustment & Search, safety, reporting, and option value \\
\bottomrule
\end{tabular}
\end{frame}

\begin{frame}{Safety as a labor constraint}
\begin{itemize}
    \item Safety changes the feasible set of jobs, not only the utility of a job already chosen.
    \item Violence and harassment can reduce applications, acceptance, hours, mobility, and retention.
    \item Reproductive autonomy changes career timing and future option value.
    \item Hidden harms can remain invisible in wages and employment status.
\end{itemize}
\end{frame}

\begin{frame}{Feasible set}
\[
\mathcal{F}_{it}
=
\{j\in\mathcal{J}_{t}:
m_{ijt}\leq \bar m_{it},
r_{ijt}\leq \bar r_{it},
q_{ijt}\geq \bar q_{it},
a_{it}\geq \bar a_{it}\}
\]
\begin{itemize}
    \item \(m\): commute or mobility burden.
    \item \(r\): safety, violence, or harassment risk.
    \item \(q\): job quality and reporting environment.
    \item \(a\): reproductive control and timing feasibility.
\end{itemize}
\end{frame}

\begin{frame}{Job utility is not just the wage}
\[
U_{ijt}=
w_{jt}
-\phi h_{jt}
-\kappa m_{ijt}
-\rho r_{ijt}
-\psi H_{ijt}
+A_{ijt}
+\Omega_{it}(T_{it})
+\varepsilon_{ijt}
\]
\begin{itemize}
    \item \(H\): hidden harm, fear, stigma, retaliation, reporting cost.
    \item \(A\): amenities, dignity, authority, job quality.
    \item \(\Omega\): option value from timing and control.
\end{itemize}
\end{frame}

\begin{frame}{Five margins to keep separate}
\begin{enumerate}
    \item Safety as a price or disutility margin.
    \item Safety as a mobility and search margin.
    \item Safety as a retention and exit margin.
    \item Safety as a hidden welfare margin.
    \item Reproductive autonomy as a timing and control margin.
\end{enumerate}
\end{frame}

\begin{frame}{Domestic violence and intimate-partner violence}
\begin{itemize}
    \item Labor mechanism: control over mobility, job search, income, and outside options.
    \item Observed margins: participation, absenteeism, hours instability, job loss, search intensity.
    \item Identification challenge: violence, work, bargaining power, and household selection are jointly determined.
    \item Useful data: police or court reports, health records, surveys, local labor shocks, employment panels.
\end{itemize}
\end{frame}

\begin{frame}{Domestic violence evidence}
\begin{itemize}
    \item Heath: labor opportunities, control over resources, and domestic violence in Bangladesh.
    \item Aizer: relative labor-market opportunities as bargaining-power variation.
    \item Frankenthal: productivity shocks and domestic violence as a newer frontier setting.
    \item Design lesson: name the direction of causality and the margin observed.
\end{itemize}
\end{frame}

\begin{frame}{Workplace violence}
\begin{itemize}
    \item Adams-Prassl, Huttunen, Nix, and Zhang link reported workplace violence to labor-market records.
    \item Observed margins: victim outcomes, coworker spillovers, separations, firm hiring, workplace composition.
    \item Identification challenge: reported incidents are selected and may reveal broader unsafe environments.
    \item Design lesson: administrative linkage can show responses ordinary wage data miss.
\end{itemize}
\end{frame}

\begin{frame}{Workplace harassment}
\begin{itemize}
    \item Folke and Rickne treat sexual harassment as a source of labor-market inequality.
    \item Harassment can generate quits, occupational segregation, and wage-risk tradeoffs.
    \item A higher wage in a hostile job is not automatically a welfare improvement.
    \item Reporting costs make complaints an equilibrium outcome, not a clean exposure measure.
\end{itemize}
\end{frame}

\begin{frame}{Power inside the workplace}
\begin{itemize}
    \item Office relationships with supervisors raise internal-labor-market questions.
    \item Margins: promotion, retention, coworker exits, perceived fairness, retaliation risk.
    \item Key challenge: separate consent, power asymmetry, favoritism, harassment risk, and spillovers.
    \item Data need hierarchy, timing, worker outcomes, and coworker outcomes.
\end{itemize}
\end{frame}

\begin{frame}{Mobility, commuting, and public-space safety}
\begin{itemize}
    \item A posted job is not accessible if the route, time, or transport mode is unsafe.
    \item Mobility risk can shrink search radius and application sets before accepted jobs appear.
    \item Safer transport can affect applications, attendance, hours, job acceptance, and time use.
    \item Non-applications are usually invisible in standard labor-force data.
\end{itemize}
\end{frame}

\begin{frame}{Mobility designs}
\begin{itemize}
    \item Randomized transport or mobility interventions directly move commute constraints.
    \item Application data reveal search radius and jobs never accepted.
    \item Time-use data reveal whether safer mobility reallocates time toward work, education, or care.
    \item Safety perceptions and local incident data help interpret the mechanism.
\end{itemize}
\end{frame}

\begin{frame}{Reproductive autonomy as labor timing}
\begin{itemize}
    \item Birth control affects expected fertility timing before pregnancy occurs.
    \item Abortion access or denial affects whether pregnancy becomes a binding work and care constraint.
    \item IVF and fertility treatment can relax biological timing constraints and provide identification.
    \item The labor object is timing, control, and option value.
\end{itemize}
\end{frame}

\begin{frame}{Birth control}
\begin{itemize}
    \item Goldin and Katz: contraception access and women's career and marriage decisions.
    \item Design family: historical and policy access variation.
    \item Observed margins: schooling, professional investment, marriage timing, career commitment.
    \item Mechanism: lower uncertainty raises returns to long human-capital investments.
\end{itemize}
\end{frame}

\begin{frame}{Abortion access and denial}
\begin{itemize}
    \item Miller, Wherry, and Foster: economic consequences of being denied an abortion.
    \item Design family: quasi-experimental comparison around provider gestational limits.
    \item Observed margins: labor-market paths, financial distress, family structure.
    \item Mechanism: control over whether and when pregnancy affects work and care.
\end{itemize}
\end{frame}

\begin{frame}{IVF and fertility treatment}
\begin{itemize}
    \item Lundborg, Plug, and Rasmussen use IVF treatment success as an instrument.
    \item Observed margins: childbearing, hours, earnings, career paths.
    \item Strength: isolates fertility timing among women seeking treatment.
    \item Limitation: local to the IVF population and treatment context.
\end{itemize}
\end{frame}

\begin{frame}{Hidden harms and reporting}
\begin{itemize}
    \item Exposure is latent: many incidents are never reported.
    \item Reporting depends on retaliation risk, stigma, documentation, job security, and trust.
    \item Labor response is selected: workers with good outside options may leave first.
    \item Welfare may improve when measured wages fall if the move is toward safety.
\end{itemize}
\end{frame}

\begin{frame}{Welfare object}
\[
V_{it}=E_t\sum_{\tau=t}^{T}\beta^{\tau-t}
\left[
c_{i\tau}
-\phi h_{i\tau}
-\rho r_{i\tau}
-\psi H_{i\tau}
+A_{i\tau}
+\Omega_{i\tau}
\right]
\]
\begin{itemize}
    \item Wages and employment are observed components, not the full welfare object.
    \item Safety, dignity, reporting costs, and option value must enter interpretation.
\end{itemize}
\end{frame}

\begin{frame}{Research designs and what they identify}
\scriptsize
\begin{tabular}{p{0.28\linewidth} p{0.31\linewidth} p{0.25\linewidth}}
\toprule
Design & What it can identify & Main threat \\
\midrule
Admin legal linkage & reported incident and labor response & selected reporting \\
Survey plus admin & latent exposure and labor outcomes & survey selection \\
Randomized transport & mobility/search constraint & spillovers and compliance \\
Legal/judicial quasi-experiment & access or denial shock & local population \\
IVF-based IV & fertility timing effect & local to IVF patients \\
Boundary discontinuity & access gradient near a rule & manipulation or sorting \\
\bottomrule
\end{tabular}
\end{frame}

\begin{frame}{Week 7 lab}
\begin{itemize}
    \item Reproduce: synthetic administrative-linkage exercise inspired by workplace violence evidence.
    \item Diagnose: classify exposure, reporting, retention, firm response, coworker spillover, and hidden welfare.
    \item Transfer: harassment-risk sorting exercise inspired by Folke and Rickne.
    \item Optional memo: office relationships, hierarchy, and workplace spillovers.
\end{itemize}
\end{frame}

\begin{frame}{Bridge to Week 8}
\begin{itemize}
    \item Safety, mobility, legal credibility, and reproductive access vary across countries and regions.
    \item The same mechanism may appear as nonparticipation, informal work, migration, segregation, or hidden welfare loss.
    \item Week 8 asks which gender-and-labor mechanisms travel and which depend on development, institutions, and norms.
\end{itemize}
\end{frame}

\begin{frame}{Takeaways}
\begin{itemize}
    \item Safety and reproductive autonomy are labor-market constraints and welfare objects.
    \item Hidden harms enter through search, mobility, sorting, quits, retention, and option value.
    \item Good evidence names the constraint, the observed margin, and the design.
    \item Wages and employment are necessary but incomplete for welfare interpretation.
\end{itemize}
\end{frame}

\end{document}
